\setcounter{page}{170}
\begin{proposition}
\begin{enumerate}
\item[(i)] For composable finite morphisms \(X\xrightarrow{f}Y\xrightarrow{g}Z\), the diagram
\[
\begin{CD}
\Rfst (gf)^b @>{\operatorname{Tr}_{gf}}>> \id \\
@VVV @| \\
\Rfst g^b \circ \Rfst f^b @>{\operatorname{Tr}_g\circ\operatorname{Tr}_f}>> \id
\end{CD}
\]
commutes as functors on \(D_{\mathrm{qc}}^+(Z)\) (resp.\ \(D_{\mathrm{qc}}(Z)\)).

\item[(ii)] Let \(u: Y' \to Y\) be a flat base extension, and \(X'=X\times_Y Y'\), \(v:X'\to X\), \(g:X'\to Y'\) as in Proposition 6.3. Then the diagram
\[
\begin{CD}
u^* \Rfst f^b @>{\sim}>> \Rfst g^b u^* \\
@V{u^*\operatorname{Tr}_f}VV @V{\operatorname{Tr}_g u^*}VV \\
u^* @= u^*
\end{CD}
\]
commutes on \(D_{\mathrm{qc}}^+(Y)\) (resp.\ \(D_{\mathrm{qc}}(Y)\)). The left vertical morphism is from [II.5.12].
\end{enumerate}
\end{proposition}

\begin{proof}
Left to the reader.
\end{proof}

\setcounter{page}{171}
\begin{theorem}[Duality for finite morphisms]
Let \(f: X \to Y\) be a finite morphism of noetherian preschemes of finite Krull dimension. Then the duality morphism
\[
\theta_{f}\colon \Rfst \RHom_X\bigl(F^\bullet,f^b G^\bullet\bigr)
\to \RHom_Y\bigl(\Rfst F^\bullet,G^\bullet\bigr)
\]
defined by composing [I.5.5] with the trace morphism \(\operatorname{Trf}_f\) is an isomorphism for all \(F^\bullet \in D_{\mathrm{qc}}^-(X)\) and \(G^\bullet \in D_{\mathrm{qc}}^+(Y)\).
\end{theorem}

\begin{proof}
Making the standard functorial reductions, we reduce to a purely algebraic statement:
Let \(A\to B\) be a ring homomorphism, \(M\) a \(B\)-module, \(N\) an \(A\)-module. Then the natural map
\[
\Hom_B\bigl(M,\Hom_A(B,N)\bigr)
\to \Hom_A(M,N)
\]
is a canonical isomorphism.
\end{proof}

\subsection{Example}
One cannot expect the duality theorem to hold for non-quasi-coherent sheaves, even for a finite étale morphism of integral noetherian schemes.

Let \(Y\) be a non-singular curve over an algebraically closed field \(k\), and \(X\to Y\) a double cover. Let \(y\in Y\) be a closed point, with two distinct points \(x_1,x_2\in X\) lying over \(y\). Let \(K(X),K(Y)\) denote the function fields of \(X,Y\) respectively.

Let \(G\) be the sheaf concentrated at \(y\) isomorphic to \(K(Y)\). Then \(G\) is an indecomposable injective \(\sheaf{O}_Y\)-module which is not quasi-coherent [II.7.11]. One verifies that \(f^b(G)\) is the direct sum of two copies of \(K(X)\), concentrated at \(x_1,x_2\) respectively; it is a direct sum of two indecomposable injective \(\sheaf{O}_X\)-modules.

\setcounter{page}{172}
Take \(F=\sheaf{O}_X\). Then:
\[
f_*\RHom_X\bigl(F,f^b G\bigr)
\]
is isomorphic to \(2\cdot K(X)\) concentrated at \(y\), while
\[
\RHom_Y\bigl(f_*F,G\bigr)
\]
is isomorphic to \(K(X)\) concentrated at \(y\).

Thus the duality morphism \(\underline{\theta}_f\) cannot be an isomorphism for this non-quasi-coherent \(G\).

\setcounter{page}{173}
\begin{remark}
Let \(f: X \to Y\) be a closed immersion of preschemes (not necessarily locally noetherian), with \(X\) a locally complete intersection in \(Y\). We may strengthen the preceding results:

We define
\[
f^{b}\colon D^{+}(Y) \to D^{+}(Y), \qquad
(\text{resp. } f^{b}\colon D(Y) \to D(X))
\]
by the same formula as before. Since \(X\) is a local complete intersection, the functor \(\RHom_Y(f_*\sheaf{O}_X,\,\cdot)\) has finite cohomological dimension.

As in Proposition 6.2, the natural morphism
\[
(g f)^{b} \to f^{b} g^{b}
\]
is defined and is an isomorphism on \(D^+(Z)\) (resp.\ \(D(Z)\)). For closed immersions \(\overline{f}^*,\overline{g}^*\) are identity functors, so no reduction to quasi-coherent complexes is needed.

The trace morphism
\[
\operatorname{Trf}_{f}\colon \Rfst f^{b}\bigl(G^\bullet\bigr) \to G^\bullet
\]
from Proposition 6.5 extends to all \(G^\bullet\in D^+(Y)\) (resp.\ \(D(Y)\)), since the morphism \(\tau\) in its construction is an identity map.

The compatibility diagrams of Proposition 6.6 remain valid. Note that for closed immersions, \(f_*\) is exact so the quasi-coherence assumption in [II.5.12] is superfluous.

Finally, the duality theorem holds for all \(F^\bullet\in D(X)\), \(G^\bullet\in D^+(Y)\). We may assume \(G^\bullet\) is a complex of injective \(\sheaf{O}_Y\)-modules; then \(f^b G^\bullet\) is also a complex of injectives, and \(f_*\) is exact. The morphism
\[
f_*\RHom_X\bigl(F^\bullet,f^b G^\bullet\bigr)
\to \RHom_Y\bigl(f_*F^\bullet,G^\bullet\bigr)
\]
is an isomorphism on each component, hence for the full complex. No appeal to the way-out functor lemma is required.

\setcounter{page}{174}
\end{remark}

\begin{proposition}
Let \(f: X \to Y\) be a finite morphism of locally noetherian preschemes.
\begin{enumerate}
\item[(a)] There is a functorial isomorphism
\[
f^b\bigl(F^\bullet \stackrel{\mathbf{L}}{\otimes} G^\bullet\bigr)
\cong f^b(F^\bullet) \stackrel{\mathbf{L}}{\otimes} f^b(G^\bullet)
\]
for \(F^\bullet \in D^{+}(Y)\), \(G^\bullet \in D^b_{\mathrm{fTd}}(Y)\).

\item[(b)] There is a functorial isomorphism
\[
\RHom\bigl(\Lfst F^\bullet,f^b G^\bullet\bigr)
\cong f^b\bigl(\RHom(F^\bullet,G^\bullet)\bigr)
\]
for \(F^\bullet \in D_{\mathrm{c}}^-(Y)\), \(G^\bullet \in D_{\mathrm{qc}}^+(Y)\).

\item[(c)] For \(F^\bullet,G^\bullet\in D_{\mathrm{qc}}^b(Y)_{\mathrm{fTd}}\), the following diagram commutes
\[
\begin{CD}
F^\bullet \otimes \Rfst f^b G^\bullet
@>{\text{Proj. Formula [II 5.6]}}>>
\Rfst\big(\Lfst F^\bullet \stackrel{\mathbf{L}}{\otimes} \Lfst G^\bullet\big)
\end{CD}
\]
\end{enumerate}
\end{proposition}

\setcounter{page}{175}
\begin{proof}
Left to the reader.
\end{proof}